\documentclass[10pt,aspectratio=1610]{beamer}
%add handout after 12pt, to disable pause
%use aspectratio=1610 to adjust 

\usepackage{amssymb}
\usepackage{amsfonts}
\usepackage{amsmath}
\usepackage{amsthm}

\usepackage{mathtools}
%\usepackage{cancel}
\usepackage[T1]{fontenc}
%\usepackage[document]{ragged2e}
%\usepackage{fancyhdr} % for header and footer
%\usepackage{tcolorbox}
%\usepackage{graphicx} %for inserting picture
%\usepackage{float} % eh centering?
%\usepackage{subfigure} % for side by side figure
%\usepackage{array} % for fixed table
\usepackage{palatino}
%\usepackage{titlesec} %for paragraph
\usepackage{indentfirst}
\usepackage{booktabs}
\usepackage{tikz-cd}

\usefonttheme{serif}
\usepackage{newpxtext,newpxmath} 

\setbeamertemplate{page number in head/foot}[totalframenumber]
\setbeamertemplate{footline}[frame number]



\usetheme{berlin}
\title{Sheaves Theory}
\author{Gilbert Allister (10123021)}
\date{\today}

\theoremstyle{definition}
%\newtheorem{definition}{Definition}

\begin{document}
    \frame{\titlepage}

    \begin{frame}{Table of Content}
        \tableofcontents
    \end{frame}

    \begin{frame}{Goal of This Presentation}
        \begin{enumerate}
            \item To understand the object of Presheaf, Stalk, Germs and Sheaf, and their morfism,
            \item To build intuition behind those objects,
            \item To understand the application of category theory.
        \end{enumerate}
        
    \end{frame}

    \section{Preliminaries}

    \begin{frame}{Some Prelim: Category Theory}
        \begin{itemize}
            \item Category and Morphism definition
            \begin{definition}[Category]
                Define $\mathcal{C}$ as category if it has \textbf{object}, \textbf{morphism},  and \textbf{composition} , such that
                \begin{itemize}
                    \item A class of $Ob\mathcal{C}$, called as class of object,
                    \item For each pair of object, a set denote $\mathcal{C}(X,Y)$ or $Hom(X, Y)$ called as morphism, from $X$ to $Y$,
                    \item For each triple $(X,Y,Z)$ objects, we have a map $\mathcal{C}(X,Y) \times \mathcal{C}(Y,Z) \to \mathcal{C}(X,Z)$ define as composition, such that it has associativity and identity, such that 
                    \begin{enumerate}
                        \item If $f \in \mathcal{C}(X,Y)$, $g \in \mathcal{C}(Y,Z)$, and $h \in \mathcal{C}(Z,W)$, we have $h(gf) = (hg)f$
                        \item For each object $X \in \mathcal{C}$, we have $id_X \in \mathcal{C}(X,X)$, such that for every $f \in \mathcal{C}(X,Y)$ and $g \in \mathcal{C}(Z,X)$, we have $f id_X = f$ and $id_X g= g$.
                    \end{enumerate}
                    
                \end{itemize}
                
            \end{definition}

            \item Examples: Sets with functions, Group with group action, Poset with $\leq$
            \item Not a category example: the category of all category(?)
        \end{itemize}
    \end{frame}

    \begin{frame}{Some Prelim: Category Theory}
        \begin{itemize}
            \item Functors
            \item Yoneda Lemma 
            \item Examples that will be used
            \begin{itemize}
                \item $\mathcal{S}ets$
                \item Ring of differentiable function $\mathcal{O}$
            \end{itemize}
            
        \end{itemize}
    \end{frame}

    \begin{frame}{Some Prelim: Category Theory}
        \begin{itemize}
            \item Colimit
            \item Adjoint
        \end{itemize}
        
    \end{frame}

    \begin{frame}{Some Prelim: Topology}
        \begin{itemize}
            \item Topological space definition
            \begin{definition}[Topology]  
                Let $X$ be a set. Define $\mathcal{T}$ as \textbf{topology}  on set $X$ if 
                \begin{enumerate}
                    \item $\emptyset$ and $X$ in $\mathcal{T}$
                    \item The union of the elements of any subcollection of $\mathcal{T}$ is in $\mathcal{T}$
                    \item The intersection of the elements of any finite subcollection of $\mathcal{T}$ is in $\mathcal{T}$
                \end{enumerate}
                We call $X$ as \textbf{topological space} if it has $\mathcal{T}$ as its topology.
                 
            \end{definition}    
            \item Open set definition
        \end{itemize}
        
    \end{frame}

    \section{Motivation}

    \begin{frame}{Motivation}
        
        \begin{itemize}
            \item Why study sheaves? 
            \item Because we need to use sheaves of function to explain schemes.
            \item Sheaf called as 'sheaf' because it glues every local data information. Presheaf is 'almost' sheaf, and it is still good idea to introduce it first.
            \item Fun Fact: Sheaves was invented by Jean Leray (a French Mathematician) in prison during WWII!
        \end{itemize}
        
    \end{frame}

    \begin{frame}{Restriction Map}
        \begin{itemize}
            \item Let $X$ be topological space. Suppose we have bunch of open sets, by let $U \subseteq X$ denote any open set and bunch of differentiable function on open set $U$. 
            \item Those differentiable functions on open set $U$ exudes a ring, denote this as $\mathcal{O}(U)$
            \item Notice if we have open set $V \subseteq U$, we can have a differentiable function at $V$. 
            \item In another word, if $V \subseteq U$ such that we have inclusion map from $V$ to $U$, we have a restriction map $res_{U, V}: \mathcal{O}(U) \to \mathcal{O}(V)$ 
             
        \end{itemize}
        
    \end{frame}

    \begin{frame}[fragile]
        \begin{itemize}
            \item Now notice if we restrict a restricted differentiable function, we still have differentiable function over there.
            \item In another word, if we have $W \hookrightarrow V \hookrightarrow U$, we can have such diagram commutes.
            \begin{figure}
                \centering
                \begin{tikzcd}
                    \mathcal{O}(U) \arrow[rrdd, "{\exists res(U, W)}"', dashed] \arrow[rrrr, "{res(U, V)}"] &  &                &  & \mathcal{O}(V) \arrow[lldd, "{res(V, W)}"] \\
                                                                                                            &  &                &  &                                            \\
                                                                                                            &  & \mathcal{O}(W) &  &                                           
                \end{tikzcd}
            \end{figure}
            
        \end{itemize}
    \end{frame}

    \begin{frame}{Sheaves Intuition}
        \begin{itemize}
            \item Now, let $U$ be 'big' open set and $\{U_i\}$ denote open covers of $U$, such that $\bigcup U_i = U$.
            \item Suppose we have $f_1$ and $f_2$ differentiable function on $U$. 
            \item If for every $U_i$, we have $f_1$ equal to $f_2$ on that $U_i$, then $f_1 = f_2$ to begin with. 
            \item In another words, $\forall \{U_i\}_{i\in I}$ and $\forall f_1, f_2 \in \mathcal{O}(U)$, we have
            \begin{align*}
                res_{U, U_i} f_1 = res_{U, U_i}f_2 \Longrightarrow f_1 = f_2
            \end{align*}
            \item This notion is called as \textbf{locality}. We will formalize this on sheaf definition.
        \end{itemize}
    \end{frame}

    \begin{frame}
        \begin{itemize}
            \item Next, with the same $U$ and $\{U_i\}$, suppose we have $f_1$ on $U_1$, $f_2$ on $U_2$ and so on, and assume they equal on the pairwise overlaps. 
            \item Notice we can just "glued" them together into one differentiable function on $U$. 
            \item Or in another words, $\forall i, f_i \in U_i$ such that $\forall i,j, res_{U_i \cap U_j} f_i = res_{U_i \cap U_j} f_j $, then 
            \begin{align*}
                \exists f \in \mathcal{O}(U) \forall i, res_{U, U_i} f = f_i
            \end{align*}
            \item This idea is called as \textbf{glueing}. Again, we will formalize this on the definition. 
            \item Notice that this will work anyways with continous function, smooth function, and just plain function. Hence, this encourages us to define such nice properties formally.
        \end{itemize}
        
    \end{frame}

    \begin{frame}{Stalk and Germ Intuition}
        \begin{itemize}
            \item Similar to previous slide's idea, can we find the intersection where the functions are equal?
            \item Then we define a relation such that 
            \begin{align*}
                (f, U)\sim(g, V) \iff \exists W \subset U \cap V, res_{U, W} \, f = res_{V, W} \, g
            \end{align*}
            for open set $U, V$ and $f \in \mathcal{O}(U)$ and $g \in \mathcal{O}(V)$
            \item One can prove this is an equivalence relation.
            \item Then we can define a class such that, 
            \begin{align*}
                \left\{(f, U): p \in U, f \in \mathcal{O}(U)\right\}
            \end{align*}
            for a point $p \in X$. 
            \item This class is known as \textbf{germs} . We will define germs using this later. 
        \end{itemize}
        
    \end{frame}

    \begin{frame}
        \begin{itemize}
            \item Consider a set of germs. Notice if we have two germs and add them together, you get another germs.
            \item Or explicitly speaking, let $f\in \mathcal{O}(U)$ and $g\in\mathcal{V}$, then $f+g$ is well-defined on $U+V$.
            \item To sketch the proof of well-defineness, suppose we have $\widetilde{f}$ has the same germs to $f$, and also $\widetilde{g}$ to $g$, 
        \end{itemize}
        
    \end{frame}


    \section{Presheaves}
    \begin{frame}{Definition}
        \begin{itemize}
            \item We will define presheaf in category theorist view.
            \begin{definition}[Presheaf]
            Define $\mathcal{F}$ as presheaf of a category $\mathcal{C}$ (can be a set, group, ring, top \dots) over a category $\mathcal{D}$ if it is a functor such that 
            \begin{align*}
                \mathcal{F}: \mathcal{D}^{op} \to \mathcal{C} (\text{or }\mathcal{S}ets, \mathcal{G}rp, \mathcal{R}ing, \mathcal{T}op)
            \end{align*}
            Or a contravariant functor from $\mathcal{D}$ to $\mathcal{C}$.
        \end{definition}
            \item For the scope of schemes theory, we let $\mathcal{D}$ be something of topological space.
        \end{itemize}
        
    \end{frame}


    \begin{frame}{Definition}
        \begin{itemize}
            \begin{definition}[Presheaf over Topological Space]
                Let $X$ be a topological space, define $\mathcal{F}$ as \textbf{presheaf over}  $X$ if it is a contravariant functor from $Open(X)$ (the category of open subsets of $X$) to a category $\mathcal{C}$. We can write it as 
                \begin{align*}
                    \mathcal{F}: Open(X)^{op} \to \mathcal{C}
                \end{align*}
            \end{definition}
            \item The category of $Open(X)$ has
            \begin{itemize}
                \item $\mathcal{O}bj$ = open subset of $U$ of $X$
                \item $\mathcal{M}or$ = inclusion $V \subseteq U$
            \end{itemize}
            \item Of course one can check whether this is a category or not. 
            
        \end{itemize}
    \end{frame}

    \begin{frame}
        \begin{itemize}
            \item This means that presheaf has following data (as a functor)
            \begin{align*}
                \forall \text{open set } U \subseteq X, U  &\mapsto \mathcal{F}(U) \text{ (as a set or group or ring \dots)} \\
                (V\subseteq U) & \mapsto (\mathcal{F}(U) \overset{res_{U, V}}{\longmapsto} \mathcal{F}(V))
            \end{align*}
            \item The identity and composition is obvious in previous explaination about restriction map.
            \item This explains why we define presheaf as contravariant.
            

        \end{itemize}
        
    \end{frame}

    \begin{frame}{Example}
        
    \end{frame}

    \section{Stalk and Germs}
    \begin{frame}
        \begin{itemize}
            \item The 'stalk' means a plant stalk, while the 'germ' means as 'cereal germ' to a stalk, as it is the heart of a function.
        \end{itemize}
        
    \end{frame}


    \section{Sheaves}

    \begin{frame}
        \begin{itemize}
            \item Locality and 
            \item Hence, presheaf failed as being 'sheaf' since it does not fulfill the locality (Uniqueness) and glueing(Existence).
        \end{itemize}
        
    \end{frame}
    
    \section{Morphism of Presheaves and Sheaves}

    \begin{frame}
        feoia
    \end{frame}

    \section{Reference}

    \begin{frame}{Reference}
        \begin{enumerate}
            \item Vakil, R. (2025). The Rising Sea: Foundations of Algebraic Geometry. Princeton University Press.
            \item Khawarizmi, F. K. (2025) Slide Seminar I for 10122040, ITB
            \item Munkres, J. R. (2014). Topology (2nd ed.). Prentice Hall, Inc.
            \item Rosiak, D. (2020). Sheaf Theory through Examples (Abridged Version). arXiv preprint arXiv:2012.08669.
        \end{enumerate}
    \end{frame}
    
    

\end{document}